% \RequirePackage{fix-cm}

\documentclass[sigconf,screen,nonacm]{acmart}        
\usepackage{graphicx}
\usepackage{natbib}
\usepackage{float}

\begin{document}

\title{A \emph{Great} Title}


\author{Antonis Protopsaltis}
\affiliation{%
  \institution{University of Western Macedonia Greece, ORamaVR Switzerland}
  \city{}
  \country{}
}

\author{Manos Kamarianakis}
\affiliation{%
  \institution{FORTH - ICS Greece, University of Crete Greece, ORamaVR Switzerland}
  \city{}
  \country{}
}


\author{George Papagiannakis}
\affiliation{%
  \institution{FORTH - ICS Greece, University of Crete Greece, ORamaVR Switzerland}
  \city{}
  \country{}
}




\begin{abstract}
This is an abstract
\keywords{keyword1 \and keyword2 \and keyword3 }
\end{abstract}



\begin{teaserfigure}
  \includegraphics[width=\textwidth, height=90pt]{example-image}
  \caption{A teaser figure.}
  \Description{A teaser figure showing an example image.}
  \label{fig:teaser}
\end{teaserfigure}

\maketitle


\section{Introduction}
This is the introduction section. 
A nice teaser figure is shown in Figure~\ref{fig:teaser}. 
A nice reference on this one is \cite{example2022}.


\section{An interesting section}
\label{sec:section1}

This is the content of the interesting section. Here is an interesting equation:
\begin{equation}\label{eq:einstein}
E = mc^2
\end{equation}

Einstein's famous equation is shown in \eqref{eq:einstein}.


\section{Another interesting section}
\label{sec:section2}

This is the content of another interesting section. In \ref{sec:section1} we discussed the previous interesting section. Now, we build upon that discussion to explore further insights, referencing Einstein's famous equation \eqref{eq:einstein}.

An important figure is shown in Figure~\ref{fig:important}. It is crucial to reference all figures appropriately in the text.

\begin{figure}[hbt]
\centering
\includegraphics[width=0.45\textwidth, height=90pt]{example-image}
\caption{An important figure. Change the 0.45 width but always less than 0.5. You can change figure to figure* if needed to have a two-column figure.}
\Description{A description of the important figure showing an example image.}
\label{fig:important}
\end{figure}





\section{Conclusions}
\label{sec:conclusions}
This is the conclusion section where we summarize the key findings and contributions of our work.

\section{Acknowledgements}
\label{sec:acknowledgements}

The work  was (partially) funded by a great EU-funded project. 
% \clearpage
\bibliographystyle{plain}
\bibliography{references} 

\end{document}


